% !TEX root = ../main.tex
\PlanPage{收支与资金安排}{武桐宇}{45}
\textbf{本页结论：}资金安排采取“先自筹验证、后以书面需求决定扩展”的顺序。下列情景只是将第 37 页的测试价格与第 44 页预算放入同一算式，不代表预测收入或融资承诺。

计算口径为：验证期净额＝（试用席位数×500 元/月×3 月）＋（专题数×5,000 元）－12,000 元预算。服务收入、固定投入与潜在数据许可必须分别记录；未收款的意向不计入收入。

\begin{table}[htbp]
\centering
\caption{验证期收支情景（均为内部假设）}
\label{tab:wty-finance-scenarios}
\small
\begin{tabularx}{\linewidth}{p{0.16\linewidth}p{0.24\linewidth}p{0.18\linewidth}X}
\toprule
情景 & 假设的已付款组合 & 假设收入 & 对资金安排的含义 \\
\midrule
审慎 & 0 个席位、0 个专题 & 0 元 & 全部 12,000 元由自筹验证预算承担；停止扩张 \\
基准 & 4 个席位连续 3 个月、1 个专题 & 11,000 元 & 仍有 1,000 元缺口；不据此宣称盈亏平衡 \\
扩展 & 8 个席位连续 3 个月、2 个专题 & 22,000 元 & 扣除验证预算后有 10,000 元余量，先用于复核与安全投入 \\
\bottomrule
\end{tabularx}
\end{table}

三个情景的共同前提是用户已付款、交付范围合规且不需新增受限数据授权；任一前提不成立即回到审慎情景。后续若需要融资，应提交实际支出、合同与合规评估，而不是以本表替代财务证明。
